\chapter{Overview}
\label{chap:overview}

Flypitch Lean 4 is a port of the original Lean 3 Flypitch development.
The long-term target is the theorem that the continuum hypothesis is
independent of ZFC. The original project reaches that theorem by combining
three layers:
\begin{itemize}
\item a first-order logic and completeness development,
\item Boolean-valued models and forcing,
\item a ZFC formalization that connects the two branches.
\end{itemize}

At the moment, the Lean 4 repository has progressed far along the
logic/completeness branch. The forcing and ZFC branch is still largely
unported. This blueprint therefore begins by documenting the material that
already exists and is checked in Lean 4, rather than presenting the full
independence proof as if it had already been reconstructed.

\section{Implemented Surface}

The verified Lean 4 surface currently includes:
\begin{itemize}
\item first-order languages, terms, formulas, and substitution,
\item a natural deduction proof system and its soundness theorem,
\item sentence and theory infrastructure,
\item compactness at the proof and theory level,
\item maximal consistent extensions,
\item directed colimits and language-extension machinery,
\item the Henkin language chain and the consistency side of henkinization.
\end{itemize}

In other words, the repository already contains the logic needed to pass from
raw first-order syntax to complete Henkin theories. The later forcing-side
modules that ultimately produce models of $\neg\mathrm{CH}$ and
$\mathrm{CH}$ are still to be ported.

\section{How To Read This Blueprint}

The chapters are organized by mathematical construction rather than by raw file
order. Each chapter explains:
\begin{itemize}
\item the mathematical object being formalized,
\item why that object is needed in the overall strategy,
\item the main Lean declarations that implement it.
\end{itemize}

The present pass starts with the first-order logic core, because every later
chapter depends on its treatment of de Bruijn variables, derivability,
satisfaction, and theories.

\section{Blueprint Roadmap}

The current blueprint is organized as a logic-side progression.

\begin{itemize}
\item Chapter~\ref{chap:fol-core} develops the first-order core: syntax,
proofs, semantics, and theories.
\item Chapter~\ref{chap:compactness-completion} explains how compactness and
completion turn consistency into complete consistent extensions.
\item Chapter~\ref{chap:colimits-language-extension} introduces the language
maps and colimits needed to compare syntax across enlarging languages.
\item Chapter~\ref{chap:henkinization} assembles those tools into the actual
henkinization construction.
\item Chapter~\ref{chap:status-next-frontier} records the present project
boundary and the next unported branch.
\end{itemize}

\section{Current Boundary}

The current frontier is not in the logic core any more. The logic-side
Henkinization machinery has been ported far enough to produce complete Henkin
extensions of consistent theories. The next genuine mathematical gap lies on
the forcing side, beginning with the still-unported \texttt{pSet\_ordinal}
development.

\section{Implemented Dependency Chain}

At the moment, the verified logic-side dependency chain can be read as:
\begin{itemize}
\item first-order syntax, proofs, semantics, and theories in
Chapter~\ref{chap:fol-core},
\item compactness and completion in
Chapter~\ref{chap:compactness-completion},
\item language colimits and extension/reflection machinery in
Chapter~\ref{chap:colimits-language-extension},
\item henkinization and complete Henkin extensions in
Chapter~\ref{chap:henkinization}.
\end{itemize}

What is still missing on the logic/completeness side is the later packaging
that would correspond to the remaining \texttt{completeness} module. Beyond
that, the next entirely unported branch begins on the forcing side at
\texttt{pSet\_ordinal}.
